\subsection{The full moving evolution and weak sources}
\label{sec:short-full-clock}

Let $C_c=(\mathbb R/c\mathbb Z)\times\mathbb R$, $c=\nu^{-2}$,
and let $v=v_c$ be the actual planar train. Fix $m>8$. Write
\[
 X_{m,c}=L^2(C_c,\langle x\rangle^{2m}dx),\qquad
 Z^{-1}_{m,c}=\{f\in(H^1_{-m,c})^*:\langle f,1\rangle=0\},
 \qquad g(s)=\exp\frac{\gamma s}{1+\nu s}.
\]
The test $1$ belongs to $H^1_{-m,c}$. The centered radial weight is
used for evolution norms. Smooth periodic weights are used only for
local smoothing, with equivalence applied at its two endpoints.

\begin{proposition}[Full forward train estimate]
\label{prop:short-full-clock}
For a fixed sufficiently large $b$ and all sufficiently small $\nu$,
the moving full complement has the exponential relative bound
\eqref{eq:short-causal-stable-moving}, the full moving evolution
satisfies \eqref{eq:short-causal-long-forward}, and on
$0<t<s\le H=\nu^{-1}((\nu b)^{-1}-1)$ its physical weak-source
extension satisfies
\[
 \left\|\frac{g(t)}{g(s)}U(s,t)\right\|_{Z^{-1}_{m,c}\to X_{m,c}}
 \le C_b\nu^{-M_s}\{1+(s-t)^{-1/2}\},\qquad
 M_s=\lceil m/2+|\Re\kappa|+2\rceil.
\]
These bounds hold on the entire zero-circulation space, with no
input parity restriction.
\end{proposition}
\begin{proof}
\emph{The full moving complement.}
Set
\[
 a=1+\nu s,\quad \theta=(\nu a)^{-1},\quad
 z=s/a=\nu^{-1}-\theta,\quad e_c=(\nu c^2)^{-1}=\nu^3,
 \quad d=c/\sqrt a=\sqrt{\theta/e_c}.
\]
The exact substitution $w(s,x)=a^{-2}W(z,x/\sqrt a)$ gives
\begin{equation}\label{eq:short-causal-moving}
 W_z=G(z)W,\qquad G=\theta^{-1}D+L_{\nu,c}(z),\qquad
 D=\Delta+\tfrac12 y\cdot\nabla+2,\qquad
 d_z/d=-1/(2\theta).
\end{equation}
The rescaled train differs from the periodized Gaussian by
$C_{j,m}e_c e^{-b_{j,m}\theta/e_c}$ in every fixed weighted derivative
norm on $\theta\ge b>0$. This follows by rescaling the train error,
interpolating its Gaussian-small unweighted norm with a larger fixed
weighted norm, and absorbing the resulting fixed powers into the
Gaussian. The required region is contained in $1\le a\le c^2$.

The common actual-cylinder comparison
\eqref{short-bg:free-cylinder}--\eqref{short-bg:interaction-cylinder}
applies in moving coordinates after including Reynolds' formula.
Here is the quantitative part needed for all parities. The paired-image
kernel satisfies, for $|x|\le d/8$, $|y|\le d/4$,
\[
 K_d(x-y)-K_\infty(x-y)
 =d^{-2}B_1(x-y)+O(d^{-4}|x-y|^3).
\]
For zero-mass input the $x$ times mass term disappears; the first
moment is retained and bounded by $C_m\|f\|_{X_m}$. The source tail
has $L^1$ norm at most $C_md^{1-m}\|f\|_{X_m}$; the separated
far kernel costs $O(d^{-1})$. The observation outside $d/8$ is
Gaussian-small. Hence the localized interaction error, also after
$d\partial_d$, is $O(d^{-2})=O(e_c/\theta)$ on the entire space.
No first-moment cancellation or parity restriction is used. The same
$O(d^{-2})$ bound holds for auxiliary nonzero-mass inputs: retain
$x\int f$ and bound it against the localized observation coefficient.
Thus restricting an ambient input to a cell creates no mass-related
gap in the compact comparison.

Let $E(z)$ be the two periodized growing vectors from
Proposition~\ref{prop:short-gaussian-seed}, and normalize their dual
rows $\Lambda(z)$ by $\Lambda E=I_2$. The planar dual kernels and
their fixed derivatives have the $O(|y|^{-2-j})$ decay obtained from
\eqref{eq:short-seed-green} and \eqref{eq:short-seed-functional}.
Paired translates therefore give dual error $O(d^{-2})$ in the
weighted dual norm, with the same bound after $d\partial_d$; direct
columns have Gaussian tails. The finite Gram inverse exists uniformly
and its differentiated error is $O(\theta^{-1}d^{-2})$.
The identities
\[
 \theta^{-1}D(P_dq)-\partial_z(P_dq)=\theta^{-1}P_dDq,
 \quad
 \theta^{-1}D^\#(P_d\ell)+\partial_z(P_d\ell)
       =\theta^{-1}P_dD^\#\ell,
 \quad D^\#=\Delta-\tfrac12y\cdot\nabla+1
\]
follow term by term using $d_z/d=-1/(2\theta)$. They are identities
for the unweighted dual kernels, with Gram normalization then applied.
In particular the combined direct and dual rows are bounded; neither
$\Lambda D$ alone nor a derivative of sharp zero extension is used.

Put $K_0=GE-E_z$. Define $H_0$ by differentiating the actual moving
pairing: $(\Lambda W)'=H_0W$ whenever $W'=GW$. Then
$M_0=H_0E=\Lambda K_0$. For $\beta_0=\gamma+\Re\kappa/\theta$ set
\begin{equation}\label{eq:short-causal-damped-complement}
 \mathcal B=G-\beta_0-EH_0+
       [-K_0+(\beta_0-1)E+EM_0]\Lambda.
\end{equation}
All added terms are bounded finite rank. Direct substitution gives
$(\Lambda W)'=-\Lambda W$ for $W'=\mathcal BW$, and
$\mathcal BE=E_z-E$. Thus $\ker\Lambda(z)$ is invariant and the
artificial pair decays at rate one. The planar limit is
$\mathscr Q(\mathscr L-\gamma)\mathscr Q-\mathscr P$, where
$\mathscr Q=I-\mathscr P$ removes both growing modes. Its full-space
semigroup is
exponentially stable by \eqref{eq:short-seed-complement}.
This statement includes the horizontal mean and every angular sector.
For ambient identifications a radial mass-one neutral vector can be
adjoined; $\mathscr L$ annihilates it and the relative damping makes
it stable. Restriction to zero circulation is preserved by the actual
equation.

To compare on a fixed Hilbert space let $I_z,R_z,\Pi_z=I_zR_z$ be
zero extension, restriction and cell projection. Extend the free
moving evolution $F$ by the actual cocycle
\[
 \overline F(z,t)=I_zF(z,t)R_t+S(z,t)(I-\Pi_t),\qquad
 Sf(y)=\rho^4f(\rho y),\quad \rho=d(t)/d(z).
\]
The two terms have disjoint entrance and exit supports. The exterior
term has norm at most $\rho^3$. For the interior, Reynolds' boundary
term cancels the dilation seam term. With $p=(1+|y|^2)^m$,
\[
 \frac d{dz}\|W\|_p^2\le
 \left[\theta^{-1}\sup\left(\frac{\Delta p}p+3
                   -\tfrac12y\cdot\nabla\log p\right)+\delta_\nu\right]
                   \|W\|_p^2.
\]
The two seam weights agree and the remaining diffusion seam term is
nonpositive. Radial Euler transport preserves the weight; the image
and actual-base velocity errors give $\sup\delta_\nu\to0$ on
$\theta\ge b$ for fixed $b$. One can take
$C_m=(m^2+3m+3)/2$ for the norm-growth contribution of $D$.
Consequently the relative free evolution satisfies
\begin{equation}\label{eq:short-causal-free-moving}
 \|\overline F_{\rm rel}(z,t)\|
 \le\exp\{[-\gamma+(C_m+|\Re\kappa|)/b+\sup\delta_\nu](z-t)\}.
\end{equation}

On each fixed block of length $B$, compact smooth tests and Duhamel
give direct strong convergence to radial Euler transport in the
successive limit $\nu\to0$, then $b\to\infty$, uniformly in the
entrance time. Density extends the convergence. The undamped norms
have limiting upper bound one, and radial transport $T$ is unitary;
the identity
\[
 \|F_n^*f-T^*f\|^2\le\|F_n\|^2\|f\|^2
                 -2\Re\langle f,F_nT^*f\rangle+\|f\|^2
\]
gives adjoint-strong convergence. The localized interaction is compact
by source and observation truncation and local Rellich compactness;
the image and dual estimates above give norm convergence of every
compact and finite-rank term in
\eqref{eq:short-causal-damped-complement}. Lemma~\ref{lem:short-compact-duhamel}
therefore compares the full-minus-free remainders in norm. If the
limiting full relative bound is $C_0e^{-hB}$, choose $B$ so that
$C_0e^{-hB}+2e^{-\gamma B}<1/4$, then $b$ large, then $\nu$ small.
Each full block contracts by $1/2$. Uniform bounds on the remaining
short block yield
\begin{equation}\label{eq:short-causal-stable-moving}
 \|S_{\rm rel}(z,t)\|\le Ce^{-h_0(z-t)},\qquad
 b\le\theta(z)\le\theta(t)\le\nu^{-1}.
\end{equation}
All block products use the same ambient radial weighted norm.

Write $W=E\alpha+r$, $\Lambda r=0$. The amplitude diagonals are
$\lambda_i+\kappa_i/\theta$; angular orthogonality makes the planar
off-diagonal $D$ pairing zero. The direct errors are bounded by
\[
 \delta_*(\theta)=C\{e_c/\theta+e_c e^{-b_2\theta/e_c}\},
\]
and both amplitude--complement couplings by
$C(\theta^{-1}+\delta_*)$. Divide each amplitude by
$e^{\lambda_i(r_0-\theta)}(r_0/\theta)^{\kappa_i}$ and the
complement by their common modulus. Equation
\eqref{eq:short-causal-stable-moving} bounds its stable convolution.
For small $\nu$, enlarge constants to bound the coupling by
$\eta(z)=C(1+e_c)/\theta(z)$ and the direct error by
$\delta(z)=Ce_c/\theta(z)$. The Gaussian error is absorbed in
these bounds. If $A_*(z)$ is the running maximum of the normalized
amplitude, the stable convolution gives
$\|r(z)\|\le Ce^{-h_0(z-t)}\|r(t)\|+C\eta(z)A_*(z)$.
The initial-complement term in the amplitude is integrable with
bound $C\|r(t)\|/b$. Equivalently, for the double convolution,
$\eta'/\eta=1/\theta\le1/b$ and $b\ge2/h_0$ give
$\int_u^z\eta(v)e^{-h_0(v-u)}dv\le(2/h_0)\eta(u)$.
Substitution into the amplitude equation and scalar Gronwall give
an exponential of
\[
 C\int_b^{1/\nu}[\delta_*(r)+(r^{-1}+\delta_*(r))^2]dr
 \le C/b+Ce_c\log(1/(\nu b))+Ce_c/b+Ce_c^2/b+o(1).
\]
The initial complement contributes only a uniformly bounded decaying
convolution. The displayed budget is uniformly bounded. We obtain,
on the entire zero-mass space and with no input symmetry restriction,
\begin{equation}\label{eq:short-causal-long-forward}
 \|U_{\rm mov}(\theta,r_0)\|_{X_{m,d(r_0)}\to X_{m,d(\theta)}}
 \le C_b e^{\gamma(r_0-\theta)}(r_0/\theta)^{\Re\kappa}.
\end{equation}

\emph{Weak sources and the polynomial loss.}
Choose the core endpoint
\[
 H=s_b=\nu^{-1}((\nu b)^{-1}-1),\qquad a_b=(\nu b)^{-1}.
\]
For $J_aW(x)=a^{-2}W(x/\sqrt a)$, change of variables gives
\begin{equation}\label{eq:short-causal-scales}
 \|J_a\|_{X_{m,d}\to X_{m,c}}\le a^{(m-3)/2},\qquad
 \|J_a^{-1}\|_{X_{m,c}\to X_{m,d}}\le a^{3/2},\qquad
 \|J_a^{-1}\|_{(H^1_{-m,c})^*\to(H^1_{-m,d})^*}\le a^{3/2}.
\end{equation}
For the last inequality,
$\langle J_a^{-1}f,\phi\rangle=a\langle f,\phi(\cdot/\sqrt a)\rangle$
and $\|\phi(\cdot/\sqrt a)\|_{H^1_{-m,c}}
\le\sqrt a\|\phi\|_{H^1_{-m,d}}$. At two endpoints
$1\le a_t\le a_s\le a_b$, the product of the first two factors is
at most $a_b^{m/2}$. Normalizing both physical cell norms by $c^{-1}$
does not change this operator estimate.

For short weak smoothing use only here the equivalent periodic weight
\[
 p_d=[1+y_2^2+(d/\pi)^2\sin^2(\pi y_1/d)]^m.
\]
Its first two logarithmic spatial derivatives are uniformly bounded,
$|\partial_z\log p_d|\le C_m/\theta$, and
$|y\cdot\nabla\log p_d|\le C_m$. The rescaled train and its first
two derivatives are uniformly bounded for $\theta\ge b$.
The adjoint localized interaction maps $X_{-m,d}$ to $H^1_{-m,d}$
uniformly: its localized source has uniformly bounded $L^1$ and
$L^2$ norms by weighted Cauchy--Schwarz; split the Biot--Savart kernel
into its integrable near part and bounded far part, and use the
uniform order-zero multiplier bound for its derivative, including
the horizontal mean. This also proves the stated bound for the
actual train perturbation.

For the reverse-time adjoint let $E_j$ be its squared weighted
$H^j_{-m}$ energy. On $\Delta z\le1$, with
$\beta=\theta^{-1}$, Reynolds cancellation and differentiated
weighted energy give
\[
 E_0'+\beta_{\min}E_1\le C_bE_0,\qquad
 E_1'+\beta_{\min}E_2\le C_b(E_1+E_0).
\]
The diffusion-weight errors have coefficient $C_m\beta$ and the
transport and compact terms have coefficient $C_b$; no negative
power of $\nu$ occurs in a Gronwall coefficient. Moreover
$\beta_{\max}/\beta_{\min}\le1+1/b$ on such an interval.
Integrate the first inequality and then the inequality for $rE_1(r)$
to get $E_1(\Delta z)\le C_b(\beta_{\min}\Delta z)^{-1}E_0(0)$.
By duality, the moving weak-to-strong norm is at most
$C_b(\beta_{\min}\Delta z)^{-1/2}$.

For $0<s-t\le1$, the clock identities are exactly
\[
 \Delta z=\frac{s-t}{a_sa_t},\qquad
 \beta_{\min}=\nu a_t,\qquad
 \beta_{\min}\Delta z=\frac{\nu(s-t)}{a_s}.
\]
Combining with \eqref{eq:short-causal-scales} gives the physical
weak-to-strong cost $C_b\nu^{-(m+2)/2}(s-t)^{-1/2}$.
For a longer gap factor at $t+1/2$ in moving coordinates and apply
\eqref{eq:short-causal-long-forward} after the short smoothing step.
The two intermediate maps $J_{a_{t+1/2}}^{-1}J_{a_{t+1/2}}$
cancel \emph{before} taking norms. Thus only one endpoint cost
$\nu^{-m/2}$ occurs, the smoothing costs at most $C_b\nu^{-1}$,
and the long-clock algebraic factor costs $\nu^{-|\Re\kappa|}$.
For the fixed integer
\[
 M_s=\left\lceil m/2+|\Re\kappa|+2\right\rceil
\]
we have on $0<t<s\le H$
\begin{equation}\label{eq:short-causal-weak-forward}
 \left\|\frac{g(t)}{g(s)}U(s,t)\right\|_{Z^{-1}_{m,c}\to X_{m,c}}
 \le C_b\nu^{-M_s}\{1+(s-t)^{-1/2}\}.
\end{equation}
Zero mass is preserved by the conservative physical equation and
then by continuous extension in the actual weighted dual source norm.

\end{proof}
